\setcounter{section}{24}

Zu einer differenzierbaren Hyperfläche

\mavergleichskette
{\vergleichskette
{ Y
}
{ \subseteq }{ \R^n
}
{  }{
}
{    }{
}
{   }{
}
}
{}{}{}
ist für jeden Punkt

\mavergleichskette
{\vergleichskette
{ P
}
{ \in }{ Y
}
{  }{
}
{    }{
}
{   }{
}
}
{}{}{}
der
\definitionsverweis {Tangentialraum}{}{}
ein Untervektorraum

\mavergleichskette
{\vergleichskette
{ T_PY
}
{ \subseteq }{ \R^n
}
{  }{
}
{    }{
}
{   }{
}
}
{}{}{}
und das Standardskalarprodukt des $\R^n$ erlaubt über die
\definitionsverweis {orthogonale Projektion}{}{}
\maabbdisp {\pi_P} {\R^n } { T_PY
} {,}
Daten im umgebenden Raum auf der Hyperfläche zu interpretieren und dadurch wichtige Konzepte für $Y$ einzuführen, wie beispielsweise die
\definitionsverweis {tangentiale Beschleunigung}{}{,}
\definitionsverweis {geodätische Kurven}{}{,}
die
\definitionsverweis {Krümmung}{}{,}
die
\definitionsverweis {kovariante Ableitung eines Vektorfeldes}{}{,}
ein
\definitionsverweis {paralleles Vektorfeld}{}{}
und den
\definitionsverweis {Paralleltransport}{}{.}
Das Standardskalarprodukt definiert aber auch eine
\definitionsverweis {riemannsche Struktur}{}{}
auf $Y$, und es erhebt sich die Frage, ob man die eben erwähnten Konzepte auch intrinsisch, nur durch Kenntnis von $Y$ und ihrer riemannschen Struktur, definieren kann
\zusatzklammer {beispielsweise für die
\definitionsverweis {hyperbolische Fläche}{}{}} {} {.}
Um diese Frage beantworten zu können, muss man eine neue Technik einführen, die sogenannten Zusammenhänge auf einem Vektorbündel und insbesondere auf dem Tangentialbündel einer
\definitionsverweis {differenzierbaren Mannigfaltigkeit}{}{.}
Im Fall einer riemannschen Mannigfaltigkeit ergibt sich in kanonischer Weise der sogenannte \stichwort {Levi-Civita-Zusammenhang} {}
\zusatzklammer {siehe
[[Riemannsche Mannigfaltigkeit/Levi-Civita-Zusammenhang/Eindeutige Existenz/Fakt|Satz 26.12]]} {} {,}
mit dem man die oben erwähnten Konzepte intrinsisch erfassen kann.

  \zwischenueberschrift{Zusammenhänge}

Es sei $X$ eine reelle
\definitionsverweis {differenzierbare Mannigfaltigkeit}{}{}
und
\maabbdisp {p} { E } { X
} {}
ein Vektorbündel über $X$, das wir ebenfalls als eine differenzierbare Mannigfaltigkeit ansetzen. Die
\definitionsverweis {Tangentialabbildung}{}{}
ist eine Abbildung
\maabb {T(p)} {TE} {TX
} {,}
die in dem kommutativen Diagramm

\mathdisp {\begin{matrix} TE & \stackrel{ T(p) }{\longrightarrow} & TX &  \\ \downarrow & & \downarrow  & \\  E  & \stackrel{ p }{\longrightarrow} &  X  & \! \! \! \!\!\!\!\!   \\ \end{matrix}} {  }

liegt, und wobei für jeden Punkt

\mavergleichskette
{\vergleichskette
{ e
}
{ \in }{ E
}
{  }{
}
{    }{
}
{   }{
}
}
{}{}{}
mit Basispunkt

\mavergleichskette
{\vergleichskette
{ x
}
{ = }{ p(e)
}
{  }{
}
{    }{
}
{   }{
}
}
{}{}{}
die
\definitionsverweis {lineare}{}{}
\definitionsverweis {Tangentialabbildung}{}{}
\maabbdisp {T_e(p)} {T_e E} { T_{x}X
} {}
vorliegt. Diese ist surjektiv, was beispielsweise daraus folgt, dass es bei einem Vektorbündel $E$ lokal durch jeden Punkt

\mavergleichskette
{\vergleichskette
{ e
}
{ \in }{ E
}
{  }{
}
{    }{
}
{   }{
}
}
{}{}{}
einen Schnitt gibt. Um eine gemeinsame Basis zu haben, fassen wir $T(p)$ als
\definitionsverweis {Bündelabbildung}{}{}
\maabbdisp {Tp} {TE} { p^* TX
} {}
von Vektorbündeln über $E$ auf. Dabei ist

\mavergleichskettedisp
{\vergleichskette
{ p^*TX
}
{ =} { E \oplus TX
}
{ =} { E \times_X TX
}
{ } {
}
{ } {
}
}
{}{}{}
die
\definitionsverweis {direkte Summe}{}{}
von Vektorbündeln über $X$, siehe
[[Topologischer Raum/Vektorbündel/Summe/Produktrealisierung/Aufgabe|Aufgabe 12.24]],
wobei wir hier $p^*TX$ als Vektorbündel über $E$ auffassen. Die Faser von $p^*TX$ über einem Punkt

\mavergleichskette
{\vergleichskette
{ e
}
{ \in }{ E
}
{  }{
}
{    }{
}
{   }{
}
}
{}{}{}
mit Basispunkt

\mavergleichskette
{\vergleichskette
{ x
}
{ = }{ p(e)
}
{  }{
}
{    }{
}
{   }{
}
}
{}{}{}
ist einfach $T_{x}X$ und die Abbildung in der Faser ist nach wie vor
\maabbdisp {T_e (p)} { T_e E } { T_{x} X
} {,}
wobei es in $p^*TX$ für jedes

\mavergleichskette
{\vergleichskette
{ e
}
{ \in }{ E
}
{  }{
}
{    }{
}
{   }{
}
}
{}{}{}
eine Kopie von $T_{x} X$ gibt.

\inputbemerkung
{}
{

Es sei
\maabb {p} { E } { X
} {}
ein differenzierbares
\definitionsverweis {Vektorbündel}{}{}
über einer
\definitionsverweis {differenzierbaren Mannigfaltigkeit}{}{}
$X$. Für eine offene Menge

\mavergleichskette
{\vergleichskette
{ U
}
{ \subseteq }{ X
}
{  }{
}
{    }{
}
{   }{
}
}
{}{}{,}
über der $E$ trivialisiert, also die Gestalt

\mavergleichskette
{\vergleichskette
{ E  \vert_U
}
{ \cong }{ U \times W
}
{  }{
}
{    }{
}
{   }{
}
}
{}{}{}
mit einem $\R$-Vektorraum $W$ hat, ist

\mavergleichskettedisp
{\vergleichskette
{ T  \left(  E \vert_U  \right)
}
{ =} { T(U \times  W )
}
{ =} { T U \times T W
}
{ =} { TU \times W \times W
}
{ } {
}
}
{}{}{,}
und die
\definitionsverweis {Tangentialabbildung}{}{}
zu
\maabb {p_U} {E \vert_U } { U
} {}
ist dabei die Projektion auf $TU$. Als Bündelhomomorphismus über $U \times W$ muss man

\mavergleichskettedisp
{\vergleichskette
{ p^*TU
}
{ =} { TU \times W
}
{ =} { TU \times_U (U \times W)
}
{ } {
}
{ } {
}
}
{}{}{}
und
\maabbdisp {} {T  \left(  E \vert_U  \right)  = TU \times W \times W  } { TU \times W
} {}
mit der Projektion auf $TU$ und der zweiten Projektion auf $W$ nehmen.

}

\inputdefinition
{}
{

Es sei
\maabb {p} { E } { X
} {}
ein
\definitionsverweis {differenzierbares Vektorbündel}{}{}
auf einer
\definitionsverweis {differenzierbaren Mannigfaltigkeit}{}{}
$X$. Das
\definitionsverweis {Kernbündel}{}{}
des
\definitionsverweis {surjektiven}{}{}
\definitionsverweis {Bündelhomomorphismus}{}{}
\maabbdisp {T(p)} {TE} {p^*TX
} {}
über $E$ heißt \definitionswort {Vertikalbündel}{.}  Es wird mit $V$ bezeichnet.

}

Das Vertikalbündel ist ein Vektorbündel über $E$ nach
[[Topologischer Raum/Reelles Vektorbündel/Homomorphismus/Kern/Untervektorbündel/Aufgabe|Aufgabe 12.19]]
und zwar ein
\definitionsverweis {Untervektorbündel}{}{}
von $TE$.

__NOEDITSECTION__

\inputfaktbeweis
{Mannigfaltigkeit/Vektorbündel/R/Vertikalbündel/Eigenschaften/Fakt}
{Lemma}
{}
{

\faktsituation {Es sei
\maabb {p} { E } { X
} {}
ein
\definitionsverweis {differenzierbares Vektorbündel}{}{}
auf einer
$C^2$-\definitionsverweis {differenzierbaren Mannigfaltigkeit}{}{}
$X$.}
\faktuebergang {Dann gelten folgende Aussagen.}
\faktfolgerung {\aufzaehlungdrei{Es liegt eine
\definitionsverweis {kurze exakte Sequenz}{}{}

\mathdisp {0 \, \stackrel{  }{\longrightarrow} \,  V   \, \stackrel{  }{ \longrightarrow}   \,  TE  \, \stackrel{  }{\longrightarrow} \, p^* TX \,\stackrel{  } {\longrightarrow}  \, 0} {  }

von
\definitionsverweis {Homomorphismen}{}{}
von Vektorbündeln über $E$ vor.
}{Für das
\definitionsverweis {Vertikalbündel}{}{}
liegt eine
\definitionsverweis {Bündelisomorphie}{}{}

\mavergleichskette
{\vergleichskette
{ V \cong p^*E
}
{ = }{ E \oplus E
}
{  }{
}
{    }{
}
{   }{
}
}
{}{}{}
vor.
}{Für jeden Punkt

\mavergleichskette
{\vergleichskette
{ e
}
{ \in }{ E
}
{  }{
}
{    }{
}
{   }{
}
}
{}{}{}
liegt eine kurze exakte Sequenz von Vektorräumen

\mathdisp {0 \, \stackrel{  }{\longrightarrow} \, E_{p(e)}   \, \stackrel{  }{ \longrightarrow}   \,  T_eE  \, \stackrel{  }{\longrightarrow} \, T_{p(e)}X \,\stackrel{  } {\longrightarrow}  \, 0} {  }

vor.
}}
\faktzusatz {}
\faktzusatz {}

}
{

\aufzaehlungzwei {Dies ist klar.
} {Es gibt einen injektiven Homomorphismus von Vektorbündeln
\maabbdisp {} {p^*E = E \oplus E} { TE
} {}
über $E$, der
\mathl{(P,u)}{,}

\mavergleichskette
{\vergleichskette
{ u
}
{ \in }{ E_P
}
{  }{
}
{    }{
}
{   }{
}
}
{}{}{,}
auf den Tangentialvektor abbildet, der durch die
\zusatzklammer {vertikale} {} {}
differenzierbare Kurve
\maabbeledisp {} {\R} {E
} { t } {(P,tu)
} {,}
gegeben ist. Unter
\maabbdisp {} {TE} { p^*TX
} {}
werden diese Tangentialvektoren auf $0$ abgebildet, da ja die repräsentierende Kurve unter
\maabb {p} { E } { X
} {}
auf $P$ abgebildet wird. Daher ist

\mavergleichskette
{\vergleichskette
{ E
}
{ \subseteq }{ V
}
{  }{
}
{    }{
}
{   }{
}
}
{}{}{}
und aus Ranggründen liegt eine Isomorphie vor.
}

}

Häufig ist das Vektorbündel $E$ gleich dem Tangentialbündel von $X$, dann liegt eine kurze exakte Sequenz

\mathdisp {0 \, \stackrel{  }{\longrightarrow} \, p^* TX  \, \stackrel{  }{ \longrightarrow}   \, TTX  \, \stackrel{  }{\longrightarrow} \, p^* TX \,\stackrel{  } {\longrightarrow}  \, 0} {  }

vor, und man muss stets sorgfältig unterscheiden, ob man das Tangentialbündel links oder rechts meint.

\inputbemerkung
{}
{

Es sei
\maabbdisp {f} {\R^n } { \R
} {}
eine zweifach
\definitionsverweis {stetig differenzierbare Funktion}{}{}
und

\mavergleichskette
{\vergleichskette
{ Y
}
{ = }{ f^{-1} (0)
}
{  }{
}
{    }{
}
{   }{
}
}
{}{}{}
die zugehörige
\definitionsverweis {differenzierbare Hyperfläche}{}{,}
die in jedem Punkt
\definitionsverweis {regulär}{}{}
sei. Das
\definitionsverweis {Tangentialbündel}{}{}
$TY$ lässt sich ebenfalls als eine Faser von Funktionen beschreiben, nämlich als

\mavergleichskettealignhandlinks
{\vergleichskettealignhandlinks
{ TY
}
{ =} { \left\{  (P,u)   \mid   f(P) = 0 , \,  \left(  D f   \right)_{ P } \left(  u  \right) = 0       \right\}
}
{ =} { \left\{  (P,u_1 , \ldots , u_n)   \mid   f(P) = 0 , \,  \sum_{i = 1}^n u_i { \frac{   \partial   f   }{  \partial   x_i   } } ( P) = 0       \right\}
}
{ \subseteq} { \R^n \times \R^n
}
{ } {
}
}
{}
{}{,}
siehe
[[Differenzierbare Hyperfläche/Tangentialbündel/Implizite Realsierung/Aufgabe|Aufgabe 10.15]].
Wir fassen dabei $f$ als Funktion auf dem $\R^{2n}$ auf, die nur von den ersten $n$ Variablen abhängt, und wir setzen

\mavergleichskettedisp
{\vergleichskette
{ g(P,u)
}
{ =} { \sum_{i = 1}^n u_i { \frac{   \partial   f   }{  \partial   x_i   } } ( P)
}
{ } {
}
{ } {
}
{ } {
}
}
{}{}{.}
Das zweite Tangentialbündel $TTY$, also das Tangentialbündel zu $TY$, können wir entsprechend mit dem
\definitionsverweis {totalen Differential}{}{}
zu
\maabbdisp {\varphi = (f,g)} {\R^{2n}} { \R^2
} {}
beschreiben, es ist

\mavergleichskettealignhandlinks
{\vergleichskettealignhandlinks
{ TTY
}
{ =} { \left\{  (P,u, v,w)  \mid   (P,u) \in TY , \,  \left(  D\varphi  \right)_{(P,u)} \left(  (v,w)  \right) = 0       \right\}
}
{ \subseteq} { \R^{4n}
}
{ } {
}
{ } {}
}
{}
{}{.}
Die
\definitionsverweis {Jacobi-Matrix}{}{,}
die das totale Differential $D \varphi$  beschreibt, ist

\mathdisp {\begin{pmatrix}   { \frac{   \partial   f   }{  \partial   x_1   } } ( P)   &   \ldots &  { \frac{   \partial   f   }{  \partial   x_n   } } ( P)  &   0 &  \ldots &  0  \\   \sum_{i = 1}^n u_i   { \frac{   \partial    }{  \partial   x_1   } } { \frac{   \partial   f   }{  \partial   x_i   } } ( P)  &  \ldots  &   \sum_{i = 1}^n u_i   { \frac{   \partial    }{  \partial   x_n   } } { \frac{   \partial   f   }{  \partial   x_i   } } ( P)   &  { \frac{   \partial   f   }{  \partial   x_1   } } ( P)  &  \ldots &  { \frac{   \partial   f   }{  \partial   x_n   } } ( P)    \end{pmatrix}} { . }

Dies ergibt neben den beiden Bedingungen für den $TY$, die sich nur auf
\mathkor {} {P} {und} {u} {}
beziehen, die beiden zusätzlichen Bedingungen an alle Variablen, nämlich

\mavergleichskettedisp
{\vergleichskette
{ \sum_{j = 1}^n v_j { \frac{   \partial   f   }{  \partial   x_j   } } ( P)
}
{ =} { 0
}
{ } {
}
{ } {
}
{ } {
}
}
{}{}{}
und

\mavergleichskettedisp
{\vergleichskette
{ \sum_{j = 1}^n v_j  \left(  \sum_{i = 1}^n u_i { \frac{   \partial    }{  \partial   x_i   } }  { \frac{   \partial   f   }{  \partial   x_j   } } ( P)  \right)      + \sum_{k = 1}^n w_k { \frac{   \partial   f   }{  \partial   x_k   } } ( P)
}
{ =} { 0
}
{ } {
}
{ } {
}
{ } {
}
}
{}{}{.}

}

\inputbemerkung
{}
{

Es sei
\maabbdisp {f} {\R^n } { \R
} {}
eine zweifach
\definitionsverweis {stetig differenzierbare Funktion}{}{}
und

\mavergleichskette
{\vergleichskette
{ Y
}
{ = }{ f^{-1} (0)
}
{  }{
}
{    }{
}
{   }{
}
}
{}{}{}
die zugehörige
\definitionsverweis {differenzierbare Hyperfläche}{}{,}
die in jedem Punkt
\definitionsverweis {regulär}{}{}
sei. Wir knüpfen an die Beschreibung des zweiten Tangentialbündels $TTY$ aus
[[Differenzierbare Hyperfläche/Zweites Tangentialbündel/Implizite Beschreibung/Bemerkung|Bemerkung 24.4]]
mit den dortigen Bezeichnungen an. Das zurückgezogene Tangentialbündel $\pi^* TY$ zu
\maabbdisp {\pi} {TY} {Y
} {,}
das zugleich das
\definitionsverweis {Vertikalbündel}{}{}
ist, ist

\mavergleichskettealign
{\vergleichskettealign
{ V
}
{ \cong} { \pi^*TY
}
{ =} { TY \times_Y TY
}
{ =} { \left\{ (P,u,z)  \mid  f(P) = 0 , \,  \sum_{i = 1}^n u_i { \frac{   \partial   f   }{  \partial   x_i   } } ( P) = 0   , \,  \sum_{k = 1}^n z_k { \frac{   \partial   f   }{  \partial   x_k   } } ( P) = 0    \right\}
}
{ } {
}
}
{}
{}{.}
Wir müssen bestimmen, wie die tangentiale Abbildung
\maabb {T(\pi)} {TTY } { \pi^* TY
} {}
und wie die Einbettung des Vertikalbündels in $TTY$ aussieht. Nach
[[Differenzierbare Hyperfläche/Tangentialbündel/Differenzierbarer Weg/Zweites Tangentialbündel/Aufgabe|Aufgabe 24.2]]
ist
\maabbeledisp {T(\pi)} {TTY} { \pi^* TY
} {(P,u,v,w)} { (P,u,v)
} {.}
Der Kern davon ist

\mavergleichskettedisp
{\vergleichskette
{ V
}
{ =} { \left\{  (P,u,0,w)   \mid   (P,u,0,w) \in TTY          \right\}
}
{ \subset} { TTY
}
{ } {
}
{ } {
}
}
{}{}{,}
wobei man direkt sieht, dass dann
\mathl{(P,u,w)}{} die Bedingung für
\mathl{\pi^*TY}{} erfüllt. Die Bündelhomomorphismen in der kurzen exakten Sequenz

\mathdisp {0 \, \stackrel{  }{\longrightarrow} \,  \pi^*TY  \, \stackrel{  }{ \longrightarrow}   \, TTY  \, \stackrel{  }{\longrightarrow} \,  \pi^*TY \,\stackrel{  } {\longrightarrow}  \, 0} {  }

kann man also in Koordinaten als

\mathl{(P,u,w) \mapsto  (P,u,0,w)}{} und
\mathl{(P,u,v,w) \mapsto (P,u,v)}{} schreiben.

}

Wenn $E$ das triviale Bündel ist, also

\mavergleichskette
{\vergleichskette
{ E
}
{ = }{ X \times \R^r
}
{  }{
}
{    }{
}
{   }{
}
}
{}{}{,}
so gibt es eine direkte Zerlegung
\zusatzklammer {über $E$} {} {}

\mavergleichskette
{\vergleichskette
{ TE
}
{ = }{ V \oplus p^* TX
}
{  }{
}
{    }{
}
{   }{
}
}
{}{}{.}
Für jeden Punkt

\mavergleichskette
{\vergleichskette
{ e
}
{ \in }{ E
}
{  }{
}
{    }{
}
{   }{
}
}
{}{}{}
repräsentieren dann die Tangentialvektoren

\mavergleichskette
{\vergleichskette
{ t
}
{ \in }{ V
}
{  }{
}
{    }{
}
{   }{
}
}
{}{}{}
die vertikalen Richtungen und

\mavergleichskette
{\vergleichskette
{ t
}
{ \in }{ p^* T_xX
}
{  }{
}
{    }{
}
{   }{
}
}
{}{}{}
die \anfuehrung{horizontalen Richtungen}{.} Für jedes Vektorbündel kann man lokal mithilfe von Trivialisierungen den Tangentialraum
\mathl{T_eE}{} in die direkte Summe aus vertikalem Raum und horizontalem Raum zerlegen. Allerdings hängen die horizontalen Unterräume wesentlich von der gewählten Trivialisierung ab und lassen sich nicht ohne Weiteres
\zusatzklammer {im Gegensatz zu den vertikalen Unterräumen, die zusammen das Vertikalbündel bilden} {} {}
zu einem Unterbündel zusammenfassen. Die Existenz eines solchen Bündels wird durch den Begriff des Zusammenhangs beschrieben.

\inputdefinition
{}
{

Es sei $X$ eine
\definitionsverweis {differenzierbare Mannigfaltigkeit}{}{}
und $E$ ein
\definitionsverweis {differenzierbares Vektorbündel}{}{}
auf $X$. Unter einem
\definitionswort {Zusammenhang}{}
auf $E$ versteht man eine direkte Summenzerlegung des
\definitionsverweis {Tangentialbündels}{}{}

\mavergleichskette
{\vergleichskette
{ TE
}
{ = }{ V \oplus H
}
{  }{
}
{    }{
}
{   }{
}
}
{}{}{}
in zwei
\definitionsverweis {Untervektorbündel}{}{}
\mathkor {} {V} {und} {H} {,}
wobei

\mavergleichskette
{\vergleichskette
{ V
}
{ \subseteq }{ TE
}
{  }{
}
{    }{
}
{   }{
}
}
{}{}{}
das
\definitionsverweis {Vertikalbündel}{}{}
ist. Das Unterbündel

\mavergleichskette
{\vergleichskette
{ H
}
{ \subseteq }{ TE
}
{  }{
}
{    }{
}
{   }{
}
}
{}{}{}
nennt man das
\definitionswort {Horizontalbündel}{}
des Zusammenhangs.

}

Ein Zusammenhang wird zumeist kurz als $\nabla$ bezeichnet, insbesondere, wenn man die durch einen Zusammenhang bestimmten Ableitungsprozesse betonen möchte. Gemäß dieser Definition ist also ein Zusammenhang nichts anderes als ein zum Vertikalbündel komplementäres Unterbündel des Tangentialbündels $TE$. Es gibt einige inhaltlich äquivalente Beschreibungen eines Zusammenhangs. Ein Zusammenhang ist äquivalent zu einem Bündelhomomorphismus
\zusatzklammer {der \stichwort {Vertikalprojektion} {} des Zusammenhangs} {} {}
\maabbdisp {\pi_{\text{vert} }} {TE} {V
} {}
mit

\mavergleichskettedisp
{\vergleichskette
{ \pi_{\text{vert} }  \circ \iota
}
{ =} {
\operatorname{Id}_{  V  }
}
{ } {
}
{ } {
}
{ } {
}
}
{}{}{,}
wobei $\iota$ die Einbettung von $V$ in $TE$ bezeichne. Das horizontale Bündel ist dann der Kern von $\pi_{\text{vert} }$, und umgekehrt liefert die Summenzerlegung eine Projektion auf die beiden Summanden. Eine solche Summenzerlegung nennt man auch eine \stichwort {Spaltung} {} der kurzen exakten Sequenz. Einen Zusammenhang kann man auf jede offene Teilmenge

\mavergleichskette
{\vergleichskette
{ U
}
{ \subseteq }{ X
}
{  }{
}
{    }{
}
{   }{
}
}
{}{}{}
einschränken.

\inputbeispiel{}
{

Auf dem trivialen Bündel
\maabbdisp {p} {X \times W = W_X } { X
} {,}
wobei $W$ einen reellen Vektorraum der
\definitionsverweis {Dimension}{}{}
$r$ und $X$ eine
\definitionsverweis {differenzierbare Mannigfaltigkeit}{}{}
bezeichnet, liegt  die kurze exakte Sequenz

\mathdisp {0 \, \stackrel{  }{\longrightarrow} \, p^*(X \times W)  =  X \times W \times W   \, \stackrel{  }{ \longrightarrow}   \, p^* T(X \times W )  = TX \times W \times W   \, \stackrel{  }{\longrightarrow} \, p^* TX  =  TX \times W   \,\stackrel{  } {\longrightarrow}  \, 0} {  }

von Vektorbündeln über
\mathl{X \times W}{} vor. Die Produkte werden dabei absolut genommen, wobei links die Identifizierung

\mavergleichskette
{\vergleichskette
{ (X \times W) \times_X (X \times W)
}
{ = }{ X \times W \times W
}
{  }{
}
{    }{
}
{   }{
}
}
{}{}{}
und rechts die Identifizierung

\mavergleichskette
{\vergleichskette
{ TX \times_X (X \times W)
}
{ = }{ TX \times W
}
{  }{
}
{    }{
}
{   }{
}
}
{}{}{}
vorgenommen wurde. Zu einem Punkt

\mavergleichskette
{\vergleichskette
{ (x,w)
}
{ \in }{ X \times W
}
{  }{
}
{    }{
}
{   }{
}
}
{}{}{}
gehört die exakte Sequenz der Fasern, also

\mathdisp {0 \, \stackrel{  }{\longrightarrow} \,  W   \, \stackrel{  }{ \longrightarrow}   \,  T_xX \times W   \, \stackrel{  }{\longrightarrow} \,  T_xX \,\stackrel{  } {\longrightarrow}  \, 0} { . }

Die triviale Spaltung dieser kurzen exakten Sequenz von Vektorbündeln nennt man den \stichwort {trivialen Zusammenhang} {} auf dem trivialen Bündel.

}

\inputbemerkung
{}
{

Es sei
\maabbdisp {f} {\R^n } { \R
} {}
eine zweifach
\definitionsverweis {stetig differenzierbare Funktion}{}{}
und

\mavergleichskette
{\vergleichskette
{ Y
}
{ = }{ f^{-1} (0)
}
{  }{
}
{    }{
}
{   }{
}
}
{}{}{}
die zugehörige
\definitionsverweis {differenzierbare Hyperfläche}{}{,}
die in jedem Punkt
\definitionsverweis {regulär}{}{}
sei. Wir knüpfen an die Beschreibung des zweiten Tangentialbündels $TTY$ aus
[[Differenzierbare Hyperfläche/Zweites Tangentialbündel/Implizite Beschreibung/Bemerkung|Bemerkung 24.4]]
und an die Bündelhomomorphismen aus
[[Differenzierbare Hyperfläche/Zweites Tangentialbündel/Induzierte Bündelhomomorphismen/Bemerkung|Bemerkung 24.5]]
an. Durch die orthogonale Projektion
\maabb {} {\R^n } { T_PY
} {}
erhält man die vertikale Projektion
\maabbeledisp {} {TTY} { V = p^*TY
} {(P,u,v,w)} { (P,u, \pi_P (w) )
} {,}
auf das
\definitionsverweis {Vertikalbündel}{}{.}
Durch die orthogonale Projektion gehört $\pi_P(w)$ zum Tangentialraum an $Y$, wenn man mit einem Vektor

\mavergleichskette
{\vergleichskette
{ (P,u,w)
}
{ \in }{ V
}
{  }{
}
{    }{
}
{   }{
}
}
{}{}{}
startet, so ist in
\mathl{(P,u,0,w)}{} das hintere $w$ automatisch im Tangentialraum, die orthogonale Projektion bildet solche Elemente also auf $(P,u,w)$ zurück ab. Es ergibt sich also ein
\definitionsverweis {Zusammenhang}{}{}
auf $TY$.

}

  \zwischenueberschrift{Vertikale Ableitung}

\inputdefinition
{}
{

Es sei $X$ eine
\definitionsverweis {differenzierbare Mannigfaltigkeit}{}{}
und $E$ ein
\definitionsverweis {differenzierbares Vektorbündel}{}{}
auf $X$, das mit einem
\definitionsverweis {Zusammenhang}{}{}
versehen sei. Unter der
\definitionswort {vertikalen Ableitung}{}
$\nabla$ versteht man die Abbildung
\maabbeledisp {\nabla} {C^1(E)} { C^0(E \otimes T^* X )
} { s } { \nabla(s) = \pi_{\text{vert} }  \circ T(s)
} {.}

}
Zu einem differenzierbaren Schnitt $s$ in $E$
\zusatzklammer {über $X$ oder einer beliebigen offenen Teilmenge

\mavergleichskettek
{\vergleichskettek
{ U
}
{ \subseteq }{ X
}
{  }{
}
{    }{
}
{   }{
}
}
{}{}{}} {} {}
wird also die Abbildung

\mathdisp {TX \stackrel{T(s)}{ \longrightarrow } TE \stackrel{ \pi_{\text{vert} } }{\longrightarrow} V \cong p^*E \longrightarrow E} {  }

zugeordnet, wobei wir

\mavergleichskette
{\vergleichskette
{ \operatorname{Hom}(TX,E)
}
{ \cong }{ E \otimes T^*X
}
{  }{
}
{    }{
}
{   }{
}
}
{}{}{}
auffassen.

Wenn zusätzlich ein
\definitionsverweis {Vektorfeld}{}{}
$F$ auf $X$, also ein Schnitt
\maabb {F} { X } { TX
} {}
im
\definitionsverweis {Tangentialbündel}{}{}
$TX$ gegeben ist, so erhält man die Abbildung
\maabbeledisp {\nabla_F} {C^1(E)} {C^0(E)
} {s } {\nabla_F(s) = \nabla (s) \circ F
} {,}
die man die \stichwort {vertikale Ableitung} {} in Richtung $F$ nennt. Man beachte, dass dabei die Abhängigkeit vom Vektorfeld $F$ nur punktweise ist, die Abhängigkeit von $s$ aber infinitesimal ist, da die Tangentialabbildung
\mathl{T(s)}{} eingeht.

__NOEDITSECTION__

\inputfaktbeweis
{Mannigfaltigkeit/Vektorbündel/R/Zusammenhang/Ableitung bezüglich Vektorfeld/Eigenschaften/Fakt}
{Lemma}
{}
{

\faktsituation {Es sei $X$ eine
\definitionsverweis {differenzierbare Mannigfaltigkeit}{}{}
und $E$ ein
\definitionsverweis {differenzierbares Vektorbündel}{}{}
auf $X$, das mit einem
\definitionsverweis {Zusammenhang}{}{}
versehen sei. Es bezeichne
\maabbdisp {\nabla_F} {C^1(E)} {C^0(E)
} {}
den zugehörigen Ableitungsoperator zu einem
\definitionsverweis {Vektorfeld}{}{}
$F$.}
\faktuebergang {Dann gelten die folgenden Eigenschaften.}
\faktfolgerung {\aufzaehlungzwei {
\mathl{\left(  \nabla_F(s)  \right)  (P)}{} hängt nur von $F(P)$
\zusatzklammer {und $s$ ab} {} {.}
} {Es ist

\mavergleichskettedisp
{\vergleichskette
{ \nabla_{g_1F_1+g_2F_2 }
}
{ =} { g_1 \nabla_{F_1 } +g_2 \nabla_{F_2 }
}
{ } {
}
{ } {
}
{ } {
}
}
{}{}{}
für Vektorfelder $F_1,F_2$ und Funktionen
\maabb {g_1,g_2} { X} { \R
} {.}
}}
\faktzusatz {}
\faktzusatz {}

}
{

\aufzaehlungzwei {Zu einem fixierten Punkt

\mavergleichskette
{\vergleichskette
{ P
}
{ \in }{ X
}
{  }{
}
{    }{
}
{   }{
}
}
{}{}{}
und einem fixierten Schnitt

\mavergleichskette
{\vergleichskette
{ s
}
{ \in }{ C^1(U,E)
}
{  }{
}
{    }{
}
{   }{
}
}
{}{}{}
zu einer offenen Menge

\mavergleichskette
{\vergleichskette
{ P
}
{ \in }{ U
}
{ \subseteq }{ X
}
{    }{
}
{   }{
}
}
{}{}{}
ist

\mavergleichskettedisp
{\vergleichskette
{ \nabla_F (s)(P)
}
{ =} { ( \pi_{\text{vert} } \circ  T_P(s) ) (F(P))
}
{ } {
}
{ } {
}
{ } {
}
}
{}{}{.}
} {Die Tangentialabbildung und die vertikale Projektion sind linear. Wegen (1) ist die Abhängigkeit von $F$ linear und auch mit der Multiplikation von Funktionen verträglich.
}

}

__NOEDITSECTION__

\inputfaktbeweis
{Differenzierbare Hyperfläche/Zweites Tangentialbündel/Induzierter Zusammenhang/Vertikale Ableitung/Fakt}
{Lemma}
{}
{

\faktsituation {Es sei

\mavergleichskette
{\vergleichskette
{ Y
}
{ \subseteq }{ W
}
{ \subseteq }{ \R^n
}
{    }{
}
{   }{
}
}
{}{}{,}
$W$
\definitionsverweis {offen}{}{,}
eine
\definitionsverweis {differenzierbare Hyperfläche}{}{}
und sei
\maabbdisp {F} { Y} { \R^n
} {}
ein tangentiales Vektorfeld.}
\faktfolgerung {Dann stimmt die
\definitionsverweis {vertikale Ableitung}{}{}
$\nabla_F$ zum
\definitionsverweis {Zusammenhang}{}{}
aus
[[Differenzierbare Hyperfläche/Zweites Tangentialbündel/Induzierter Zusammenhang/Bemerkung|Bemerkung 24.8]]
mit der
\definitionsverweis {kovarianten Ableitung}{}{}
überein.}
\faktzusatz {}
\faktzusatz {}

}
{

Die vertikale Ableitung längs $F$  des Zusammenhangs eines differenzierbaren Schnittes
\maabb {s} { Y } { TY
} {}
ist

\mathdisp {Y \stackrel{F}{\longrightarrow} TY \stackrel{T(s)}{\longrightarrow} TTY  \stackrel{  \pi_{\text{vert} } }{\longrightarrow} TY} { . }

Dabei bildet $T(s)$ einen Punkt
\mathl{(P,u)}{} auf
\mathl{(P,u, s(P), \left(  D s   \right)_{ P } \left(  u   \right))}{}  in der Notation von
[[Differenzierbare Hyperfläche/Zweites Tangentialbündel/Induzierter Zusammenhang/Bemerkung|Bemerkung 24.8]]
ab und dieses wird auf
\mathl{(P, \pi_P  \left(  \left(  D s   \right)_{ P } \left(  u   \right)  \right) )}{} abgebildet. Dies stimmt mit der Definition der kovarianten Ableitung überein.

}

 [[Kategorie:Latexseite]]
